\medskip

Une concession automobile vend des véhicules à moteur électrique et des véhicules à moteur thermique.

Certains clients, avant de se rendre sur le site de la concession, ont consulté la plate- forme numérique de la concession. On a ainsi observé que :
\begin{itemize}
\item[$\bullet$~~]20\,\% des clients sont intéressés par les véhicules à moteur électrique et 80\,\% préfèrent s'orienter vers l'achat d'un véhicule à moteur thermique ;
\item[$\bullet$~~]lorsqu'un client souhaite acheter un véhicule à moteur électrique, la probabilité pour que le client ait consulté la plate-forme numérique est de $0,5$ ;
\item[$\bullet$~~]lorsqu'un client souhaite acheter un véhicule à moteur thermique, la probabilité pour que le client ait consulté la plate-forme numérique est de $0,375$.
\end{itemize}

\begin{list}{\textbullet}{On considère les évènements suivants:}
\item $C$ : \og un client a consulté la plate-forme numérique \fg{} ;
\item $E$ : \og un client souhaite acquérir un véhicule à moteur électrique \fg{} ;
\item $T$ : \og un client souhaite acquérir un véhicule à moteur thermique \fg.
\end{list}

Les clients font des choix indépendants les uns des autres.

\medskip

\begin{enumerate}
\item 
	\begin{enumerate}
		\item On résume les informations dans un arbre pondéré.

\begin{center}
\bigskip
  \pstree[treemode=R,nodesepA=0pt,nodesepB=4pt,levelsep=3.5cm,nrot=:U]{\TR{}}
 {
 	\pstree[nodesepA=4pt]{\TR{$E$}\naput{$0,2$}}
 	  {
 		\TR{$C$}\naput{$0,5$}
 		\TR{$\overline{C}$}\nbput{\blue $1-0,5=0,5$}	   
 	  }
 	\pstree[nodesepA=4pt]{\TR{$T$}\nbput{$0,8$}}
 	  {
 		\TR{$C$}\naput{$0,375$}
        \TR{$\overline{C}$}\nbput{\blue $1-0,375=0,625$} 
     }
}
\bigskip
\end{center}
		
La probabilité qu'un client choisi au hasard souhaite acquérir un véhicule à moteur électrique et ait consulté la plate-forme numérique est:

$P(E\cap C)= P(E)\times P_{E}(C)=0,2\times 0,5 = 0,1$.
		

		\item 
D'après la formule des probabilités totales:

$P(C)=P(E\cap C) + P(T\cap C)= 0,1+0,8\times 0,375 = 0,4$.
		
		\item On suppose qu'un client a consulté la plate-forme numérique; la probabilité qu'il souhaite acheter un véhicule à moteur électrique est :

$P_{C}(E)= \dfrac{P(E\cap C)}{P(C)}=\dfrac{0,1}{0,4} = 0,25$.
	\end{enumerate}
\item La concession accueille quotidiennement $17$ clients en moyenne.
On note $X$ la variable aléatoire donnant le nombre de clients souhaitant acquérir un véhicule à moteur électrique.

	\begin{enumerate}
		\item 
Les clients font des choix indépendants les uns des autres, donc on a une répétition de 17 épreuves indépendantes dont la probabilité de succès de chacune est $p=0,2$.

La variable aléatoire $X$ suit donc la loi binomiale de paramètres $n=17$ et $p=0,2$.

		\item La probabilité qu'au moins trois des clients souhaitent acheter un véhicule à moteur électrique lors d'une journée est $P(X\geqslant 3)$.
		
$P(X\geqslant 3) = 1-P(X<3) = 1- P(X\leqslant 2) \approx 1- \np{0,3096}$

Donc la probabilité cherchée est d'environ $0,69$.

	\end{enumerate}
\end{enumerate}

\bigskip
